\documentclass[11pt,letterpaper]{article}
\usepackage[T1]{fontenc}
\usepackage[utf8]{inputenc}
\usepackage{lmodern,geometry,microtype,xcolor,enumitem,hyperref,fancyhdr}
\geometry{margin=0.78in}
\definecolor{ink}{HTML}{17324D}
\definecolor{accent}{HTML}{126A78}
\hypersetup{colorlinks=true,linkcolor=accent,urlcolor=accent,breaklinks=true}
\setlength{\parindent}{0pt}
\setlength{\emergencystretch}{3em}
\setlist[itemize]{leftmargin=*,itemsep=3pt,topsep=3pt}
\pagestyle{fancy}\fancyhf{}\fancyhead[L]{\small Architecture Documentation Suite}\fancyhead[R]{\small v1.0}\fancyfoot[R]{\thepage}\setlength{\headheight}{15pt}
\begin{document}
\thispagestyle{empty}
{\color{accent}\large TECHNICAL REFERENCE \hfill NOVEMBER 2026}\par\vspace{1.1cm}
{\color{ink}\LARGE\bfseries\raggedright Software Architecture Document (SAD)}\par\vspace{0.35cm}
{\Large Reusable Views and Beyond architecture-description template}\par\vspace{0.55cm}
\rule{\textwidth}{0.7pt}\par\smallskip
{\small Architecture Documentation Suite | Release 1.0}\par\vspace{0.7cm}
\tableofcontents
\newpage
\section*{Document control}\addcontentsline{toc}{section}{Document control}
[Fill in: entity/system of interest, organization, document owner, version, date, status, baseline commit or release, reviewers, change history, change-request channel, distribution and classification.] Clearly designate current-state versus proposed-state content. Remove instructional placeholders only after supplying real evidence or a justified not-applicable entry.\par\smallskip
\section*{1. Documentation roadmap}\addcontentsline{toc}{section}{1. Documentation roadmap}
[State purpose, scope, exclusions, intended readers, related documents, how the views are organized and how each stakeholder can navigate to the material needed.] Record the expected reader tasks and the acceptance decision for this architecture description (AD).\par\smallskip
\section*{2. System overview}\addcontentsline{toc}{section}{2. System overview}
[Describe mission, external context, principal capabilities, constraints, lifecycle, architecture drivers, critical dependencies and high-level quality attributes. State what is unknown.] Include a system context presentation with explicit boundaries and legend. Identify the baseline environment and source provenance.\par\smallskip
\section*{3. Stakeholders, concerns and requirements}\addcontentsline{toc}{section}{3. Stakeholders, concerns and requirements}
[Create one row per concern: C-ID; stakeholder role; concern; priority; measurable question; view IDs; requirement IDs; evidence; status.] Record high-priority quality attribute scenarios with source, stimulus, environment, affected artifact, response and response measure. Distinguish target from measured result.\par\smallskip
\section*{4. Viewpoint and model-kind definitions}\addcontentsline{toc}{section}{4. Viewpoint and model-kind definitions}
[For each viewpoint: VP-ID; purpose; stakeholders/concerns; applicable model kinds; elements, relations, properties and constraints; notation; construction method; analysis method; consistency checks.] Define identifiers, relation directions, external/internal elements and treatment of missing information. [Describe how each model kind will be represented and interpreted.]\par\smallskip
\section*{5. Selected views and view packets}\addcontentsline{toc}{section}{5. Selected views and view packets}
[Maintain a view register: V-ID, viewpoint, title, scope, concerns, model sources, owner, status, packet links and omission rationale.] Repeat the following template per selected view or view packet. Do not mechanically include every view if it serves no concern.\par\smallskip
\subsection*{5.x.1 Primary presentation}
[Insert diagram or table, figure ID, scope, version, notation key, meaningful element names and typed relationship labels. Explain omitted detail and link to related packets.]\par\smallskip
\subsection*{5.x.2 Element catalog}
[Record element ID, name, type, responsibility, properties, owner, evidence and status. Record relation ID, source, target, type, direction, protocol/contract, conditions, evidence and status. For each interface specify provided/required role, syntax, semantics, security, timing, failure modes, retries, compatibility and contract tests. Document relevant behavioral traces or lifecycle state models.]\par\smallskip
\subsection*{5.x.3 Context}
[Represent external dependencies and neighboring units using this view's vocabulary; link to broader context and scope. Distinguish assumptions about environment from known connections.]\par\smallskip
\subsection*{5.x.4 Variability guide}
[For each variation point: VPt-ID; alternative options; default; binding time; selection mechanism; constraints; security implications; verification. Write not applicable with explanation if no supported variation is known.]\par\smallskip
\subsection*{5.x.5 Rationale}
[Reference ADRs; explain style selection, assumptions, alternatives, tradeoffs, and how the view addresses the identified concerns. If rationale is unavailable from implementation evidence, mark it unknown.]\par\smallskip
\section*{6. Beyond views: correspondence and traceability}\addcontentsline{toc}{section}{6. Beyond views: correspondence and traceability}
[Record typed mappings between module, runtime, deployment and other views. Include one-to-many/many-to-one mappings, correspondence rules, known inconsistencies and their dispositions. Link concerns to views, elements, decisions, requirements and tests. Record unaddressed concerns explicitly.]\par\smallskip
\section*{7. Architecture decisions and crosscutting rationale}\addcontentsline{toc}{section}{7. Architecture decisions and crosscutting rationale}
[ADR-ID, title, date, status, owner, problem, drivers, alternatives, decision, rationale, positive/negative consequences, dependencies, validation, supersession and review trigger.] Separate observed design from proposed improvements. Record security, reliability, scalability, cost, maintainability, privacy and data governance tradeoffs as relevant.\par\smallskip
\section*{8. Risk, assumptions and open issues}\addcontentsline{toc}{section}{8. Risk, assumptions and open issues}
[Risk ID, description, likelihood, impact, affected elements and concerns, existing controls, mitigation, owner, evidence, due date and acceptance decision.] [Assumption ID, source, confidence, validation plan, owner.] [Issue ID, inconsistency or gap, disposition, owner.] Never imply an unknown control exists.\par\smallskip
\section*{9. Verification and release}\addcontentsline{toc}{section}{9. Verification and release}
[Record evidence register, review question sets, reviewers, results, outstanding findings, exceptions and approval. Include build status, diagram verification, correspondence checks and source baseline.] Define triggers for revision: interface changes, new trust boundary, deployment topology change, decision supersession or stakeholder need. Link to document 05 for measurable local gates.\par\smallskip
\section*{10. Directory}\addcontentsline{toc}{section}{10. Directory}
[Add element index, view/packet directory, terms, acronyms, source register and related artifacts.] Source note: view and beyond-view organization adapted from Clements et al. chapter 10; viewpoint and traceability additions are local extensions informed by the ISO 42010 public overview. This template is a starting point, not a conformance determination.\par\smallskip
\section*{Bibliographic foundation}
Paul Clements et al., \emph{Documenting Software Architectures: Views and Beyond}, 2nd ed., Addison-Wesley/Pearson, copyright 2011, ISBN 978-0-321-55268-6. Section and chapter references refer to the supplied edition. Additional framework sources are identified in the applicable documents.\par
\end{document}
